% ============================================================================
% CAPÍTULO 1 — INTRODUÇÃO
% ============================================================================
\chapter{Introdução}
\label{cap:introducao}

A teoria financeira tradicional parte do pressuposto de que investidores
tomam decisões de forma racional, buscando maximizar sua utilidade
esperada a partir das informações disponíveis. As finanças comportamentais
questionam esse pressuposto ao mostrar que fatores psicológicos, como
heurísticas, vieses cognitivos e diferenças individuais de personalidade,
influenciam sistematicamente as decisões financeiras, afastando-as do
ideal puramente racional \cite{jiangpengyan2024}. Dentro dessa perspectiva,
um instrumento amplamente utilizado pelo mercado financeiro para orientar
recomendações de investimento é o \textbf{perfil de investidor}: um
levantamento do quanto um indivíduo está disposto e é capaz de tolerar
risco, tipicamente classificado em categorias como conservador, moderado e
arrojado. No Brasil, a verificação desse perfil (\textit{suitability}) é
inclusive uma exigência regulatória para instituições que prestam serviços
de investimento \cite{cvm30_2021}.

Ainda que o perfil de investidor seja operacionalizado, em geral, por meio
de questionários padronizados sobre horizonte de tempo, patrimônio e
tolerância declarada ao risco, uma linha de pesquisa em finanças
comportamentais busca relacionar esse perfil a características de
personalidade mais estáveis do indivíduo, em particular ao modelo
\textbf{Big Five} de personalidade (Capítulo \ref{cap:referencial}).
Jiang, Peng e Yan \cite{jiangpengyan2024}, em uma pesquisa com 3.325
investidores reais, encontraram evidências de que traços de
personalidade, destacadamente Abertura à Experiência e Neuroticismo,
se relacionam com crenças sobre o mercado, aversão ao risco e alocação de
carteira, com poder explicativo comparável ao de variáveis demográficas
tradicionais. De Bortoli et al. \cite{debortoli2019}, em um estudo com
124 estudantes universitários, verificaram que, entre os cinco traços,
apenas Abertura à Experiência se associou ao perfil de risco observado em
uma simulação de investimento. Esses dois estudos servem de base teórica e metodológica
para o presente trabalho e são discutidos em detalhe no Capítulo
\ref{cap:referencial}.

Paralelamente, o avanço recente de modelos de linguagem de grande escala
(\textit{Large Language Models}, doravante LLMs) tem aberto uma nova possibilidade
metodológica para as ciências comportamentais: usar agentes de
inteligência artificial como respondentes simulados em pesquisas que, com
sujeitos humanos, demandariam tempo e recursos consideráveis. Um LLM pode
ser instruído, por meio de um \textit{prompt} de sistema, a responder como
se tivesse determinadas características, como um perfil de
personalidade específico, o que permite gerar, em escala e a baixo
custo, um número arbitrariamente grande de "agentes" com personalidade
controlada experimentalmente. Essa abordagem carrega limitações de
validade externa evidentes (um agente de IA não é um investidor real, e
seu comportamento reflete os dados e vieses do modelo que o gera), mas
oferece, em contrapartida, um ambiente de simulação totalmente controlado:
o nível de personalidade "verdadeiro" de cada agente é conhecido por
construção, sem o erro de medida inerente a autorrelatos humanos.

É nesse cruzamento entre a literatura de personalidade e perfil de
investidor, de um lado, e o uso de LLMs como sujeitos simulados, de outro,
que se insere este trabalho: agentes de IA recebem uma personalidade
Big Five induzida por instrução, respondem a um questionário de
autorrelato de personalidade (para verificar se a personalidade induzida
foi de fato "incorporada" pelo agente) e a um questionário de perfil de
investidor, e a relação entre o nível de personalidade induzido e o perfil
de investidor obtido é analisada estatisticamente, em busca de padrões
comparáveis aos encontrados por Jiang, Peng e Yan \cite{jiangpengyan2024} e
De Bortoli et al. \cite{debortoli2019}.

\section{Problema de pesquisa}

A pergunta de pesquisa que orienta este trabalho é: \textbf{níveis
induzidos de personalidade Big Five em agentes de inteligência artificial
predizem o escore do questionário de perfil de investidor?}

Essa pergunta geral é desdobrada em três questões de pesquisa (RQ, do
inglês \textit{research questions}), respondidas ao longo do trabalho e
retomadas de forma explícita no Capítulo \ref{cap:conclusao}:

\begin{description}
    \item[RQ1:] Os LLMs conseguem simular consistentemente os traços de
    personalidade solicitados?
    \item[RQ2:] Quais associações entre personalidade induzida e perfil de
    investidor são robustas?
    \item[RQ3:] Quanto essas associações e os perfis financeiros dependem
    do modelo e dos parâmetros de geração?
\end{description}

A RQ1 trata da validade da própria simulação: sem uma indução de
personalidade verificada, nenhuma associação com o perfil de investidor
pode ser interpretada. A RQ2 separa, entre as associações observadas na
análise principal, as que se repetem sob diferentes condições das que
dependem de uma condição específica. A RQ3 avalia em que medida os
resultados são uma propriedade do modelo de linguagem e da configuração
de geração usados, e não da relação entre personalidade e risco em si.

Investigar essas questões contribui, por um lado, para avaliar se relações
já documentadas entre personalidade e comportamento de risco financeiro em
seres humanos se reproduzem em um ambiente de simulação com agentes de IA,
o que teria implicações tanto para a validade desse tipo de simulação
em pesquisas comportamentais quanto para a compreensão de possíveis vieses
desses modelos ao lidar com temas financeiros. Por outro lado, o desenho
experimental adotado (indução controlada de personalidade, seguida de
checagem de manipulação) oferece um paralelo metodológico direto ao
desenho de De Bortoli et al. \cite{debortoli2019}, permitindo comparação
relativamente direta entre os achados deste trabalho e os da literatura de
referência.

\section{Justificativa}

Este trabalho se justifica em três frentes complementares. Em primeiro
lugar, do ponto de vista teórico, ele testa, em um ambiente controlado e de
baixo custo, se relações já estabelecidas na literatura de finanças
comportamentais entre personalidade e comportamento de investimento se
replicam em um contexto de simulação com agentes de IA: gerar 300
respondentes simulados como os deste trabalho seria inviável em tempo e
custo com sujeitos humanos reais.

Em segundo lugar, do ponto de vista metodológico, o trabalho contribui
para a discussão, ainda incipiente, sobre o uso de LLMs como sujeitos
simulados em pesquisa comportamental: ao comparar explicitamente diferentes
modelos de linguagem quanto à qualidade da indução de personalidade (seção
\ref{sec:selecao-modelo}), o trabalho evidencia que essa qualidade não é
uniforme entre modelos e que a checagem de manipulação (verificar se o
agente de fato incorporou a personalidade pedida antes de interpretar
qualquer efeito sobre a variável dependente) é uma etapa metodológica
necessária, não apenas desejável, nesse tipo de simulação.

Em terceiro lugar, do ponto de vista prático, entender melhor os
determinantes psicológicos do perfil de investidor é relevante para a
adequação de produtos financeiros ao perfil do cliente (\textit{suitability})
e para iniciativas de educação financeira, ainda que as conclusões deste
trabalho, por partirem de agentes simulados, devam ser interpretadas como
exploratórias em relação a esse uso prático (ver limitações, seção
\ref{sec:limitacoes-pipeline}).

\section{Objetivos}

\subsection{Objetivo geral}

Investigar se traços de personalidade do modelo Big Five, induzidos por
instrução em agentes de inteligência artificial baseados em modelos de
linguagem de grande escala, predizem o perfil de investidor desses
agentes, medido por um questionário de perfil de investidor (IPA).

\subsection{Objetivos específicos}

\begin{enumerate}
    \item Induzir, por meio de instrução de sistema (\textit{prompt}),
    níveis (Baixo/Médio/Alto) dos cinco traços do Big Five em agentes
    baseados em LLM, e verificar a fidelidade dessa indução por meio de um
    questionário de autorrelato de personalidade (checagem de
    manipulação);
    \item Aplicar um questionário de perfil de investidor (IPA) aos
    mesmos agentes e obter, para cada um, um escore contínuo e uma
    classificação categórica (conservador, moderado ou arrojado);
    \item Testar a associação estatística (correlação de Spearman e
    regressão OLS) entre o nível de personalidade induzido e o escore do
    IPA, para cada um dos cinco traços do Big Five;
    \item Comparar os achados obtidos com os resultados relatados por
    De Bortoli et al. \cite{debortoli2019} e por Jiang, Peng e Yan
    \cite{jiangpengyan2024}, identificando convergências e divergências;
    \item Comparar o comportamento de diferentes modelos de linguagem
    (\texttt{llama3.2:1b}, \texttt{llama3.2:3b}, \texttt{gemma3:4b} e
    \texttt{llama3.1:8b}) quanto à qualidade da indução de personalidade,
    de modo a identificar em que medida o tamanho/família do modelo afeta a
    validade da simulação;
    \item Testar se induzir os 5 traços do Big Five simultaneamente no
    mesmo agente, em comparação a induzir apenas 1 traço por vez, afeta a
    qualidade da indução de personalidade ou a força da associação com o
    escore do IPA;
    \item Avaliar a robustez e a generalização das associações encontradas,
    repetindo a análise principal com outras temperaturas de geração, com
    um novo sorteio de perfis de personalidade e com outros modelos de
    linguagem.
\end{enumerate}

\section{Etapas do trabalho: seleção de modelo, piloto e análise principal}
\label{sec:trabalho-realizado-intro}

Como etapa preliminar à análise apresentada neste trabalho, foram testados
quatro modelos de linguagem distintos (\texttt{llama3.2:1b},
\texttt{llama3.2:3b}, \texttt{gemma3:4b} e \texttt{llama3.1:8b}) quanto à
sua capacidade de incorporar de forma consistente a personalidade induzida
por instrução (seção \ref{sec:selecao-modelo}). Esse levantamento mostrou
que a qualidade da indução de personalidade varia substancialmente entre
modelos: o modelo mais leve testado (\texttt{llama3.2:1b}) não incorporou
nenhum dos cinco traços de forma estatisticamente detectável, ao passo que
o modelo finalmente adotado como base para a análise principal deste
trabalho (\texttt{llama3.1:8b}) apresentou manipulação forte e
significativa em todos os cinco traços. Esse resultado é, em si, um
achado metodológico relevante: ele indica que o tamanho e a família do
modelo de linguagem são pré-requisitos para a validade desse tipo de
simulação, e que a checagem de manipulação não pode ser tratada como uma
formalidade dispensável.

O dataset descrito acima, baseado no instrumento de autorrelato BFI-10,
constitui o \textbf{estudo piloto} deste trabalho (seção
\ref{sec:bfi10-piloto}). Foi esse piloto que motivou tanto a escolha do
\texttt{llama3.1:8b} quanto o primeiro indício do achado central deste
trabalho, a associação entre Abertura à Experiência e o escore do IPA.

A partir desse piloto, a \textbf{análise principal} deste trabalho migrou
para um instrumento de autorrelato mais granular, o BFI-44 (44 itens,
seção \ref{sec:exp-a-b}), com dois desenhos de indução testados: o
\textbf{Experimento A}, que repete o desenho do piloto (os 5 traços
induzidos simultaneamente em cada agente) e constitui os resultados
principais reportados neste trabalho (Capítulo \ref{cap:resultados}); e
o \textbf{Experimento B}, um teste de robustez metodológica em que cada
agente tem apenas 1 traço induzido por vez, motivado pela pergunta de se
a indução simultânea de personalidade dilui o sinal de cada traço
individual. Os resultados do Experimento A confirmam e fortalecem o
achado de Abertura à Experiência observado no piloto, e revelam efeitos
adicionais de Extroversão, Consciência e Neuroticismo.

Por fim, uma \textbf{análise de robustez e generalização} repete o
Experimento A variando, uma de cada vez, a temperatura de geração do
modelo, o sorteio dos perfis de personalidade e o próprio modelo de
linguagem (três modelos abertos adicionais com N=300 e dois modelos
maiores, acessados por API, em caráter exploratório). Essa análise
mostra que os efeitos de Abertura à Experiência e de Extroversão se
mantêm em todas as condições testadas, ao passo que os de Consciência e
Neuroticismo mudam de significância e até de sinal conforme a condição,
e por isso são tratados neste trabalho como \textbf{resultados
instáveis} (Capítulos \ref{cap:resultados} e \ref{cap:discussao}).

\section{Estrutura do trabalho}

Além desta introdução, este trabalho está organizado em mais cinco
capítulos. O Capítulo \ref{cap:referencial} apresenta o referencial
teórico, cobrindo o conceito de personalidade e o modelo Big Five, o
conceito de perfil de investidor em finanças comportamentais, uma
revisão detalhada dos dois artigos-base que fundamentam as hipóteses deste
trabalho e os estudos recentes sobre personalidade induzida e risco em
LLMs. O Capítulo \ref{cap:metodologia} descreve o desenho experimental,
os instrumentos utilizados, o pipeline de geração dos agentes, o processo
de seleção do modelo de linguagem e os métodos de análise estatística. O
Capítulo \ref{cap:resultados} apresenta os resultados obtidos: checagem
de manipulação, distribuição do perfil de investidor, associações entre
personalidade induzida e escore do IPA e a análise de robustez e
generalização. O Capítulo \ref{cap:discussao}
interpreta esses resultados à luz da literatura de referência e discute as
limitações do estudo. Por fim, o Capítulo \ref{cap:conclusao} responde às
três questões de pesquisa, retoma os objetivos específicos, sintetiza as contribuições
do trabalho e aponta direções para pesquisas futuras.
